%% VI. Thảo luận: diễn giải theo câu hỏi nghiên cứu, hàm ý, hạn chế. Được access_vi.tex \input.
\section{Thảo luận}
\label{sec:discussion}

\subsection{Phân tích thiết kế đã đúng ở đâu (RQ1)}
Bài đồng hành hứa rằng thông tin Fisher chính xác của M1 có thể dùng để lên kế hoạch một nghiên cứu lấy mẫu trải nghiệm. Trên hai thiết kế thật, với lịch hỏi không đều, một phần ba epoch thiếu và các phiên bị loại bởi quy tắc nhiễu giả, cận tính trước khi khoá, sai số chuẩn Fisher tại ước lượng và sai số bootstrap của $\kinh$ chênh nhau nhiều nhất 8\,\% trên ds001787 và 32\,\% trên OSF 3uszv. Trên ds001787, 9,9\,\% đạt được nằm giữa cận tính trước khi khoá và cận tính lại tại tỉ số tín hiệu trên nhiễu đo được (9,1--13,3\,\%); trên OSF 3uszv, 8,5\,\% đạt được vượt cận tính lại (7,4\,\%) 15\,\%. Với $\kd$, cận tính trước khi khoá lạc quan 1,8 lần và 1,2 lần. Thứ tự giữa hai bộ cũng được dự đoán: hỏi dày hơn và phiên dài hơn của OSF 3uszv cho cận nhỏ hơn và sai số đạt được nhỏ hơn. Một phân tích thiết kế kiểu này vì thế dùng được cho mục đích của nó, quyết định cần bao nhiêu người, phiên dài bao nhiêu và hỏi bao nhiêu lần, miễn là sai số nó nêu được hiểu là sai số Fisher của một mô hình được giả định là đúng.

\subsection{Điều nó không lường trước được (RQ1)}
Một cận không lường trước được sai đặc tả, và kênh công suất alpha sai đặc tả trên cả hai bộ: sai số theo thông tin quan sát và theo profile của $\kinh$, $\sA$ và $\rA$ vượt sai số Fisher 2,1--2,5 lần trên ds001787 và 1,1--1,4 lần trên OSF 3uszv, và phần vượt biến mất trên dữ liệu mô phỏng từ fit. Phép chẩn đoán này rẻ, vì chỉ cần một Hessian và một bộ dữ liệu mô phỏng, và chúng tôi khuyến nghị dùng nó bên cạnh bootstrap, vốn không nhìn thấy sai đặc tả vì mô phỏng từ chính mô hình. Trong thực hành, các con số của bản đồ thiết kế nên được đọc như cận dưới mỗi khi một đại diện chưa chắc tuân theo mô hình, và khoảng dựa trên hàm hợp lý nên được báo cáo bên cạnh khoảng Fisher; trên ds001787 hai tiêu chí cho H1 cho kết luận khác nhau.

\subsection{Mô hình quan sát là nút thắt (RQ2)}
Hai đại diện chúng tôi đăng ký không mang thông tin phát hiện được về báo cáo ngoài phụ thuộc chuỗi của chính báo cáo, ở bộ nào cũng vậy. Kênh độ phức tạp là nhiễu trên ds001787 và, trên OSF 3uszv, là một tín hiệu 40~s không liên quan đến báo cáo; alpha vùng sau không gắn với điểm ở bộ nào. Động lực của M1 xác định được và được ước lượng với độ chính xác đã dự đoán, nhưng hàm đọc ra từ đặc trưng EEG sang trạng thái được báo cáo mới là nơi thiếu thông tin. Điều này phù hợp với quan hệ gián tiếp giữa alpha vùng sau và DMN mà các nghiên cứu ghi đồng thời EEG--fMRI tìm thấy \cite{Laufs2003,Mantini2007}. Tìm một đặc trưng EEG bám theo trạng thái được báo cáo, với tỉ số tín hiệu trên nhiễu trong dải mà phân tích thiết kế cho thấy là cần thiết, là bài toán mở mà lớp mô hình này để lại cho người làm thực nghiệm, và các bộ dữ liệu lấy mẫu trải nghiệm hiện có \cite{Weinstein2018,Brandmeyer2018,RodriguezLarios2021} là nơi để tìm.

\subsection{Hàm ý cho khung BQI (RQ3)}
Không bộ nào ủng hộ các dự đoán T1, T2, T4 và T5 của BQI, còn T3 chỉ đúng về dấu (Mục~\ref{sec:res-bqi}), nhưng phép so sánh không cùng loại. Các giá trị BQI được fit từ tóm tắt phân tích gộp của tương phản fMRI và tương quan cắt ngang \cite[Mục~7.2--7.5]{BQI}, một hạn chế mà chính BQI nêu \cite[Mục~11.3]{BQI}, còn các ước lượng ở đây đến từ đại diện EEG quan sát theo thang phút trong từng phiên. Một hằng số thời gian fit từ tóm tắt nhóm của phân tích gộp cùng lắm mô tả cách các tương phản DMN ở mức nhóm biến đổi qua quá trình thực hành; còn 0,76 và 7,1~phút ước lượng ở đây mô tả cách công suất alpha vùng sau dao động quanh trung bình của phiên. Vì DMN Ornstein--Uhlenbeck của M1 là trường hợp thế bậc hai của động lực Langevin trong BQI \cite[pt.~(83)--(87)]{BQI}, $\kinh$ ước lượng được là độ cong của thế đó cho đại diện EEG, không phải cho chính DMN. Hệ số ghép âm trên OSF 3uszv nói rằng, với alpha vùng sau đứng thay cho DMN, alpha và độ phức tạp tăng cùng nhau, ngược với hệ số ghép trong \cite[pt.~(53)]{BQI}; điều đó không bác bỏ việc pt.~(53) đúng với chính DMN. Điều phép so sánh cho thấy là bước mà BQI còn để ngỏ: trước khi ước lượng được các tham số của nó từ EEG, các đặc trưng EEG phải được hiệu chuẩn theo biến trạng thái, chẳng hạn bằng các bản ghi đồng thời EEG, fMRI và thước đo độ phức tạp mà chính BQI đề xuất \cite[Mục~14]{BQI}, và QoC phải được thao tác hoá bằng một công cụ đã được kiểm chứng. Các phần cấu trúc của BQI mà M1 giữ lại, một quá trình DMN Ornstein--Uhlenbeck nuôi một biến tích hợp và một hàm đọc ra tuyến tính sang báo cáo, đều ước lượng được trên cả hai bộ với độ chính xác mà phân tích thiết kế đã dự đoán.

\subsection{Hằng số tốc độ là tính chất của đại diện (RQ4)}
Phép lặp lại tách hai loại kết quả. Các kết quả về phương pháp đứng vững giữa các phòng thí nghiệm: sai số dự đoán và đạt được của $\kinh$ khớp nhau trong một phần ba, kênh alpha có cùng dấu hiệu, EEG không thêm gì cho dự đoán, và dị biệt giữa người tập rất lớn. Các kết quả về nội dung thì không: hằng số thời gian DMN chuyển từ 0,8 sang 7~phút, hằng số thời gian tích hợp từ 5 sang 0,6~phút, và khác biệt nhóm về $\kinh$ mà ds001787 cho thấy đã biến mất. Mô phỏng ở Mục~\ref{sec:simmethods} cho thấy điều này là đương nhiên chứ không khó hiểu. Fit Ornstein--Uhlenbeck trả về một thoả hiệp có trọng số phương sai giữa các thang thời gian có trong đặc trưng; công suất alpha vùng sau trong một phiên mantra 45~phút và trong một phiên hơi thở 60~phút, ghi với montage khác nhau, không nhất thiết chứa cùng hỗn hợp. Bài học không phải là một trong hai ước lượng sai, mà là một hằng số tốc độ ước lượng từ đại diện chưa hiệu chuẩn là tính chất của đại diện đó trong bản ghi đó. Các tham số mà một phân tích thiết kế hứa hẹn mang lại chỉ có ý nghĩa ngang với ánh xạ từ đặc trưng sang biến trạng thái, và ánh xạ đó phải được thiết lập và chứng tỏ là ổn định giữa các phòng thí nghiệm trước khi dùng bản đồ thiết kế để lên kế hoạch cho một nghiên cứu hay một can thiệp vòng kín \cite{Sitaram2017,Ros2020}: các hệ số điều khiển tính trong bài đồng hành giả định tham số có nghĩa.

\subsection{Đăng ký trước đã bắt được gì}
Hai tình tiết minh hoạ giá trị của khoá và cổng. Quy tắc nhiễu giả đã khoá của ds001787 loại 45\,\% số phiên, điều tiêu chí dừng cho phép; không có khoá, phản ứng tự nhiên sẽ là chuyển sang quy tắc thích ứng, vốn cho hằng số thời gian tích hợp khác và $\beta$ dương, rồi báo cáo kết quả đó. Trên OSF 3uszv, một biến thể độ nhạy chép từ bản đăng ký thứ nhất không chạy được, và 17 phiên lúc đầu được tiền xử lý trên một máy không có gói ICLabel, điều mà log tiền xử lý đã ghi lại; chúng được xử lý lại trước khi báo cáo bất kỳ fit nào, và cả hai việc được liệt kê ở Phụ lục~\ref{app:deviations}. Đăng ký trước không làm phân tích đơn giản hơn, nhưng làm cho danh sách những gì được quyết định sau khi xem dữ liệu trở nên ngắn và rõ ràng \cite{Nosek2018}.

\subsection{Hạn chế}
\label{sec:limitations}
Các phép thử thực nghiệm là hai bộ dữ liệu có hằng số tốc độ lệch nhau tám đến chín lần. Trên ds001787, quy tắc nhiễu giả đã đăng ký loại 18 trong 40 phiên một cách không ngẫu nhiên, và so sánh nhóm bị nhiễu bởi giới và tuổi; trên OSF 3uszv, riêng cửa sổ hỏi đã bỏ một phần ba EEG, điểm tự chấm hai cực và độc lập theo chuỗi, và mỗi phiên có một đoạn đầu trước lần hỏi đầu tiên mà việc che nó làm hằng số thời gian DMN tăng gấp đôi. Ở cả hai, đặc trưng EEG là đại diện chưa có ánh xạ hiệu chuẩn sang $A$ hay $\Phi$ (alpha vùng sau không phải thước đo hoạt động DMN, độ phức tạp Lempel--Ziv không phải thước đo $\Phi$), các mức $A_0$, $\Phi^*$ và $\gamma_0$ không ước lượng được sau khi quy tâm, và kênh alpha có dấu hiệu sai đặc tả. Hai thang điểm tự chấm là các đại lượng thao tác đứng thay cho QoC của BQI, chưa phải thước đo đã được kiểm chứng của nó, và đại diện độ phức tạp, cũng như chỉ số của BQI, đo độ phức tạp chứ không phải tích hợp \cite[Mục~11.3]{BQI}. M1 tuyến tính với nhiễu Gauss và hàm đọc ra tuyến tính, trong khi điểm tự chấm là thứ bậc và bị chặn, hai cực ở một bộ, và động lực thiền có thể có các chuyển trạng thái rời rạc \cite{Hasenkamp2012,Christoff2016}; các mở rộng phi tuyến sẽ cần bộ lọc xấp xỉ. Các cá nhân dùng chung một vectơ tham số trong fit gộp, dù độ lệch chuẩn giữa người của $\log\kinh$ là 1,2--1,5 đơn vị log; mô hình phân cấp là bước tiếp theo tự nhiên và sẽ thay đổi bản đồ thiết kế về lượng. Cuối cùng, phân tích được lên kế hoạch, khoá và chạy bởi một tác giả, và các bản đăng ký trước không được phản biện trước khi khoá.
